\section{Sweep Results and Design-Region Screening}

The first result layer establishes the baseline behavioural response and anchors the later sweeps. It is a reproducible response case generated by the MATLAB behavioural model, not a complete design validation (Figure~\ref{fig:baseline-response}).

The feedback-capacitance sweep shows the central compensation trade-off. In the selected behavioural sweep, increasing $C_f$ reduces extracted bandwidth and generally reduces peaking. The bandwidth and peaking figures are interpreted together because each captures only one side of the design trade-off (Figures~\ref{fig:cf-bandwidth} and~\ref{fig:cf-peaking}). Compensation choices are therefore screened jointly for bandwidth and peaking rather than selected from bandwidth alone.

The design-region map extends the feedback-capacitance sweep across a broader parameter space. It shows how selected $C_{pd}$ and $f_t$ assumptions affect project-defined Safe, Marginal, or Risky outcomes (Figure~\ref{fig:design-region-map}; Table~\ref{tab:screening-criteria}). The representative response figure provides example response shapes for selected categories (Figure~\ref{fig:representative-responses}).

Together, these sweep results provide the behavioural screening layer of the report. They make the compensation and finite-bandwidth trends visible before any claim is made about agreement with vendor macromodels.
